\documentclass[10pt,twoside,letterpaper]{article}
\usepackage[utf8]{inputenc}
\usepackage{amsmath,amsfonts,amssymb}
\usepackage[spanish,es-noshorthands]{babel}
\usepackage[T1]{fontenc}
\usepackage{lmodern}
\usepackage{graphicx,hyperref}
\usepackage{tikz,pgf}
\usepackage{marvosym}
\usepackage{printsudoku}
\usepackage{multicol}
\usepackage[total={7.5in,9.75in},centering]{geometry}
\usepackage{fancyhdr}
\pagestyle{fancy}
\fancyhead[LE]{Colegio Arborizadora Baja}
\fancyhead[RE]{PEI:``Hacia una cultura para el desarrollo sostenible''}
\fancyfoot[RO]{\Email gavendanor@colarborizadorabaja.edu.co}
\fancyhead[LO]{\url{www.autistici.org/mathgerman}}
\fancyfoot[RE]{NIT 830024976-8, Código DANE: 11100103084-8}
\fancyfoot[LE]{Calle 59I \#44A - 02 \Telefon 7313994 - 7313995}
\fancyhead[RO]{Nit 830024976-8, Código DANE 11100103084-8}

\author{Germ\'an Avenda\~no Ram\'irez~\thanks{Lic. Mat. U.D., M.Sc. U.N.}}
\title{\begin{minipage}{.2\textwidth}
\includegraphics[height=1.75cm]{Images/logo-colegio.png}\end{minipage}
\begin{minipage}{.55\textwidth}
\begin{center}
Animaplano 09\\
Matemáticas $7^{\circ}$
\end{center}
\end{minipage}\hfill
\begin{minipage}{.2\textwidth}
\includegraphics[height=1.75cm]{Images/logo-sed.png} 
\end{minipage}}
\date{}
\usetikzlibrary{arrows}
\newcommand{\degre}{\ensuremath{^\circ}}
\definecolor{qqwuqq}{rgb}{0,0.39,0}
\begin{document}
\maketitle
Nombre: \hrulefill Curso: \underline{\hspace{1cm}}  Fecha: \underline{\hspace{2cm}}
\paragraph*{Temas:} Operaciones aritméticas básicas, potenciación, operaciones entre conjuntos y elementos básicos de geometría.
\paragraph*{Aprendizajes esperados:} Repaso de las operaciones aritméticas básicas y de algunas operaciones básicas entre conjuntos.
\subsection*{Cuestionario}
Conteste las preguntas del cuestionario, haciendo los procedimientos necesarios al frente de cada pregunta. Las respuestas numéricas se ubican en el plano dispuesto para ello. Como ayuda, se han puesto notas de pie de páginas.

También como reto se ha puesto un sudoku fácil.
\begin{enumerate}
  \item Sea $ A=\{2,17,9\} $ $ B=\{2,5,8,17\} $. El producto de los elementos de $ A\cap B= $\\ \textbf{Nota:} $A\cap B$~\footnote{\emph{Recuerde que $A\cap B$ está formado por los elementos que pertenecen a $A$ y $B$ simultáneamente}}
 \item El triple del número 11:
  \item Halle $ (9^5\div9^3)+2^1= $
  \item Halle $ (10^6\div10^4)-4^1= $
  \item Halle $ 10^2 $, menos $ (8\times2!)$ \textit{Recuerde como se define $n!$}~\footnote{Recuerde que el factorial de un número natural $ n! $ es $ n!=n\times(n-1)\times(n-2)\times\ldots\times3\times2\times1 $, es decir el número por todos sus antecesores.} 
  \item $ 1000-947= $
  \item Si $A+B=37$, entonces $(A+10)+(B+27)=$
  \item El 43\% de 200
  \item 1 centena menos 1 docena 
  \item Sume al triple del número 9, el cuádruple del número 10:
  \item[11] $ (1000\div10)-(7\times3)= $
  \item[12] $ (-7\times-7)+(5\times2)= $
  \item[13] $ (-5\times2)+48= $
  \item[14] $ (-1\times-2\times-5)+36= $
  \item[15] El triple del número 13
  \item[16] La tercera parte de 180
  \item[17] Si $p+14=43$, entonces $p=$
  \item[18] Sume 4 veces, el número 6
  \item[19] Si el perímetro~\footnote{El perímetro de una figura es la suma de sus lados} de un cuadrado mide 16. ¿Cuál es su área?
  \item[20] Halle $ 4!-\sqrt{25}= $
  \item[21] Sea $ A=\{15,5,9\} $ $ B=\{7,15,9,5\} $. Halle la raíz cuadrada de la suma de los elementos de $ A\cup B= $
  \item[22] $A-B$ en grados
  \usetikzlibrary{arrows}
\definecolor{qqwuqq}{rgb}{0,0.39,0}
\begin{tikzpicture}[line cap=round,line join=round,>=triangle 45,x=1.0cm,y=1.0cm,scale=1.5]
\draw[->,color=black] (-1.02,0) -- (1.46,0);
\filldraw[color=qqwuqq,fill=qqwuqq,fill opacity=0.1] (-.35,0) rectangle (0,.35);
\foreach \x in {-1,1}
\draw[shift={(\x,0)},color=black] (0pt,-2pt);
\draw[->,color=black] (0,-0.44) -- (0,1.54);
\foreach \y in {,1}
\draw[shift={(0,\y)},color=black] (2pt,0pt) -- (-2pt,0pt);
\clip(-1.02,-0.44) rectangle (1.46,1.54);
\draw [shift={(0,0)},color=qqwuqq,fill=qqwuqq,fill opacity=0.1] (0,0) -- (0:0.8) arc (0:51:0.8) -- cycle;
\draw [shift={(0,0)},color=qqwuqq,fill=qqwuqq,fill opacity=0.1] (0,0) -- (51:0.6) arc (51:90:0.6) -- cycle;
\draw [domain=-1.02:1.46] plot(\x,{(-0--1.23*\x)/1});
\draw (-0.74,0.64) node[anchor=north west] {$A$};
\draw (0.06,0.94) node[anchor=north west] {$B$};
\begin{scriptsize}
\draw (1.02,0.44) node {$51\textrm{\degre}$};
\end{scriptsize}
\end{tikzpicture}
  \item[23] 1 siglo menos 18 años
  \item[24] 1/2 siglo + 2 años
  \item[25] La tercera parte de 99
  \item[26] 1 decena, más 1 docena
  \item[27] El número de vértices de un polígno de 12 lados
  \item[28] El número de lados de un polígono con 13 vértices
  \item[29] Halle $ 96\div4 $
  \item[30] En años, 6 lustros más 4 años
  \item[31] Si $ m+28=95 $, entonces $ m= $
\end{enumerate}

  \begin{minipage}{.5\textwidth}
\begin{center}
\begin{tikzpicture}[scale=.9]
 \fill (1,0) node[above]{1} circle (0.2ex);
 \fill (2,0) node[above]{2} circle (0.2ex);
 \fill (3,0) node[above]{3} circle (0.2ex);
 \fill (4,0) node[above]{4} circle (0.2ex);
 \fill (5,0) node[above]{5} circle (0.2ex);
 \fill (6,0) node[above]{6} circle (0.2ex);
 \fill (7,0) node[above]{7} circle (0.2ex);
 \fill (8,0) node[above]{8} circle (0.2ex);
 \fill (9,0) node[above]{9} circle (0.2ex);
 \fill (10,0) node[above]{10} circle (0.2ex);
 \fill (1,-1) node[left]{11} circle (0.2ex);
 \fill (2,-1) circle (0.2ex);
 \fill (3,-1) circle (0.2ex);
 \fill (4,-1) circle (0.2ex);
 \fill (5,-1) circle (0.2ex);
 \fill (6,-1) circle (0.2ex);
 \fill (7,-1) circle (0.2ex);
 \fill (8,-1) circle (0.2ex);
 \fill (9,-1) circle (0.2ex);
 \fill (10,-1) circle (0.2ex);
 \fill (1,-2) node[left]{21} circle (0.2ex);
 \fill (2,-2) circle (0.2ex);
 \fill (3,-2) circle (0.2ex);
 \fill (4,-2) circle (0.2ex);
 \fill (5,-2) circle (0.2ex);
 \fill (6,-2) circle (0.2ex);
 \fill (7,-2) circle (0.2ex);
 \fill (8,-2) circle (0.2ex);
 \fill (9,-2) circle (0.2ex);
 \fill (10,-2) circle (0.2ex);
 \fill (1,-3) node[left]{31} circle (0.2ex);
 \fill (2,-3) circle (0.2ex);
 \fill (3,-3) circle (0.2ex);
 \fill (4,-3) circle (0.2ex);
 \fill (5,-3) circle (0.2ex);
 \fill (6,-3) circle (0.2ex);
 \fill (7,-3) circle (0.2ex);
 \fill (8,-3) circle (0.2ex);
 \fill (9,-3) circle (0.2ex);
 \fill (10,-3) circle (0.2ex);
 \fill (1,-4) node[left]{41} circle (0.2ex);
 \fill (2,-4) circle (0.2ex);
 \fill (3,-4) circle (0.2ex);
 \fill (4,-4) circle (0.2ex);
 \fill (5,-4) circle (0.2ex);
 \fill (6,-4) circle (0.2ex);
 \fill (7,-4) circle (0.2ex);
 \fill (8,-4) circle (0.2ex);
 \fill (9,-4) circle (0.2ex);
 \fill (10,-4) node[right]{50} circle (0.2ex);
 \fill (1,-5) node[left]{51} circle (0.2ex);
 \fill (2,-5) circle (0.2ex);
 \fill (3,-5) circle (0.2ex);
 \fill (4,-5) circle (0.2ex);
 \fill (5,-5) circle (0.2ex);
 \fill (6,-5) circle (0.2ex);
 \fill (7,-5) circle (0.2ex);
 \fill (8,-5) circle (0.2ex);
 \fill (9,-5) circle (0.2ex);
 \fill (10,-5) circle (0.2ex);
 \fill (1,-6) node[left]{61} circle (0.2ex);
 \fill (2,-6) circle (0.2ex);
 \fill (3,-6) circle (0.2ex);
 \fill (4,-6) circle (0.2ex);
 \fill (5,-6) circle (0.2ex);
 \fill (6,-6) circle (0.2ex);
 \fill (7,-6) circle (0.2ex);
 \fill (8,-6) circle (0.2ex);
 \fill (9,-6) circle (0.2ex);
 \fill (10,-6) circle (0.2ex);
 \fill (1,-7) node[left]{71} circle (0.2ex);
 \fill (2,-7) circle (0.2ex);
 \fill (3,-7) circle (0.2ex);
 \fill (4,-7) circle (0.2ex);
 \fill (5,-7) circle (0.2ex);
 \fill (6,-7) circle (0.2ex);
 \fill (7,-7) circle (0.2ex);
 \fill (8,-7) circle (0.2ex);
 \fill (9,-7) circle (0.2ex);
 \fill (10,-7) circle (0.2ex);
 \fill (1,-8) node[left]{81} circle (0.2ex);
 \fill (2,-8) circle (0.2ex);
 \fill (3,-8) circle (0.2ex);
 \fill (4,-8) circle (0.2ex);
 \fill (5,-8) circle (0.2ex);
 \fill (6,-8) circle (0.2ex);
 \fill (7,-8) circle (0.2ex);
 \fill (8,-8) circle (0.2ex);
 \fill (9,-8) circle (0.2ex);
 \fill (10,-8) circle (0.2ex);
 \fill (1,-9) node[left]{91} circle (0.2ex);
 \fill (2,-9) circle (0.2ex);
 \fill (3,-9) circle (0.2ex);
 \fill (4,-9) circle (0.2ex);
 \fill (5,-9) circle (0.2ex);
 \fill (6,-9) circle (0.2ex);
 \fill (7,-9) circle (0.2ex);
 \fill (8,-9) circle (0.2ex);
 \fill (9,-9) circle (0.2ex);
 \fill (10,-9) node[right]{100} circle (0.2ex);
\end{tikzpicture}
\end{center}
\end{minipage}\hfill
\cluefont{\Large}
\cellsize{1.75\baselineskip}
\begin{minipage}{0.45\linewidth}\begin{center}
SE5 (easy) \\
\sudoku{se5.sud}
\end{center}\end{minipage}
\end{document}
